\chapter{Классификация существующих скалярных носителей центрального источника}
\label{ch:v8-baryon-c0-existing-scalar-source-carrier-classification}

\section{Контракт повторного использования}

Предыдущий гейт допустил условный одномерный носитель $s$, но не
идентифицировал его внутри уже построенной теории. Буквальное повторное
использование требует одновременно
\begin{equation}
 \mathcal C_s=(\dim_{\mathbb R}s=1,\ s\in\mathbb R,\ Gs=s,\
 SO(3)_{\rm fam}s=s,\ [\Gamma,s]=0,\ s\in\mathcal F_{\rm active},\
 \partial_s\partial_\lambda S\ne0).
 \label{eq:v8-existing-scalar-literal-contract}
\end{equation}
Последнее условие существенно: совпадения квантовых чисел недостаточно,
если общий родитель не содержит билинейного портала к
$h^0=\lambda Q$.

\section{Хиггсовское направление}

Активное поле $H\sim(\mathbf1,\mathbf2)_{1/2}$ не является
калибровочным синглетом:
\begin{equation}
 J_2\binom10=\binom0{-1}\ne0,
 \qquad
 J_2=\begin{pmatrix}0&1\\-1&0\end{pmatrix}.
 \label{eq:v8-existing-scalar-higgs-nonsinglet-witness}
\end{equation}
Его квадратичный инвариант
\begin{equation}
 I_H=H^\dagger H\ge0,
 \qquad
 \chi_H=I_H-\frac{v^2}{2}\ge-\frac{v^2}{2}
 \label{eq:v8-existing-scalar-higgs-centered-halfline}
\end{equation}
вещественен, gauge- и family-инвариантен, но является составной
полупрямой, а не свободной знаковой координатой $\mathbb R$.
Симметрия разрешает портал
\begin{equation}
 V_{H Q}=-\kappa_H\lambda I_H,
 \label{eq:v8-existing-scalar-higgs-q-portal}
\end{equation}
однако минимальное direct-sum завершение Тома IV давало
\begin{equation}
 \kappa_H^{\rm inherited}=0,
 \qquad
 v\ \text{и}\ \kappa_H\ \text{не выведены слепо}.
 \label{eq:v8-existing-scalar-higgs-parent-boundary}
\end{equation}
Следовательно, $I_H$ является допустимой формой источника, но не уже
полученным источником.

\section{Радиус когерентности стрелок}

В Томе VII условно выведен rank-one конденсат
\begin{equation}
 B_*=\sqrt3\,uv^*,
 \qquad \|u\|=\|v\|=1,
 \qquad \rank B_*=1.
 \label{eq:v8-existing-scalar-coherence-condensate}
\end{equation}
Его радиальный инвариант имеет точную нормировку
\begin{equation}
 I_B=\Tr(BB^*)\ge0,
 \qquad
 I_B(B_*)=3.
 \label{eq:v8-existing-scalar-coherence-radius}
\end{equation}
Поэтому $I_B$ уже несёт ненулевое gauge-, family- и grading-инвариантное
математическое ожидание. Центрированная координата
\begin{equation}
 \chi_B=I_B-3\ge-3
 \label{eq:v8-existing-scalar-coherence-centered-domain}
\end{equation}
снова не является полной вещественной линией, но для роли составного
источника это не обязательно. Разрешённый смешанный член равен
\begin{equation}
 V_{BQ}=-\kappa_B\lambda I_B,
 \qquad
 j_{\rm eff}=\kappa_B I_B(B_*)=3\kappa_B.
 \label{eq:v8-existing-scalar-coherence-q-portal}
\end{equation}
В наследованном действии Тома VIII такой межносительный член отсутствует:
\begin{equation}
 \left.
 \frac{\partial^2S_{\rm inherited}}
 {\partial\lambda\,\partial I_B}
 \right|_{B_*}=0.
 \label{eq:v8-existing-scalar-coherence-portal-absent}
\end{equation}
Итак, $I_B$ является лучшим существующим носителем значения, но не
происхождения коэффициента $\kappa_B$.

\section{Нескалярные и неактивные кандидаты}

Гладкий относительный объект предыдущей части имеет
\begin{equation}
 \rank_{\mathbb C}B(A)=6,
 \qquad B(A_0)=0,
 \label{eq:v8-existing-scalar-smooth-order-type}
\end{equation}
то есть он операторнозначен и не даёт ненулевой одномерный источник на
физическом фоне. Проекторный параметр порядка и семейный носитель имеют
типы
\begin{equation}
 Q_{\rm proj}\in\operatorname{Sym}^2_0(\mathbb R^3),
 \quad \dim_{\mathbb R}=5,
 \qquad
 \Sigma\in\mathbb R^3,
 \label{eq:v8-existing-scalar-family-nonsinglet-types}
\end{equation}
и оба не являются семейными синглетами. Это видно уже по свидетелям
\begin{equation}
 [L_{12},\operatorname{diag}(1,-1,0)]\ne0,
 \qquad
 L_{12}(1,0,0)^T\ne0.
 \label{eq:v8-existing-scalar-family-action-witnesses}
\end{equation}

Разность двух вещественных копий twisted-ветви и NCG-singlet стерильного
сектора имеют правильный буквальный тип. Но первая требует неактивного
удвоения алгебры, динамически закрытого для обычного следа, а второй
отсутствует в действующем $H_{15}$. Их статус можно записать как
\begin{equation}
 a(\sigma_{\rm tw})=a(\sigma_{\nu_R})=0,
 \qquad
 a(I_H)=a(I_B)=1,
 \label{eq:v8-existing-scalar-activity-indicators}
\end{equation}
где $a=1$ означает принадлежность активному носителю.

\section{Итог классификации}

Из восьми проверенных классов ни один не выполняет буквальный контракт:
\begin{equation}
 N_{\rm literal}=0/8.
 \label{eq:v8-existing-scalar-literal-ledger}
\end{equation}
После ослабления требования «фундаментальная знаковая линия» до
«составной скалярный источник» остаются ровно два кандидата,
\begin{equation}
 \mathcal S_{\rm composite}=\{I_H,I_B\},
 \qquad
 N_{\rm composite}=2/2,
 \qquad
 N_{\rm inherited\ portal}=0/2.
 \label{eq:v8-existing-scalar-composite-ledger}
\end{equation}
Только $I_B$ уже имеет точное ненулевое значение, полученное внутри
условного родительского действия. Поэтому классификация выбирает его как
следующий объект проверки, но не объявляет физическую щель выведенной.

\begin{theorem}[Уникальный существующий составной кандидат]
В активной базе нет буквального одномерного знакового Real-singlet с уже
унаследованной связью к $Q$. Среди активных составных скалярных инвариантов
симметрий остаются $I_H$ и $I_B$, причём только радиус когерентности
$I_B=\Tr(BB^*)$ имеет условно выведенное ненулевое вакуумное значение
$I_B(B_*)=3$.
\end{theorem}

\begin{nogocriterion}[Инвариант не равен порталу]
Нельзя отождествлять наличие ненулевого $I_B$ с происхождением источника
$j$. До вывода единого родителя с ненулевым смешанным гессианом
$\partial_\lambda\partial_{I_B}S$ коэффициент $\kappa_B$, знак щели и
физическая нормировка остаются новыми данными.
\end{nogocriterion}

\begin{protocol}\begin{enumerate}
\item проверенных классов носителей: \textbf{$8$};
\item буквальных активных повторных использований: \textbf{$0$};
\item активных составных singlet-инвариантов: \textbf{$2$};
\item Higgs-норма имеет выведенный масштаб и портал: \textbf{нет};
\item радиус когерентности имеет точный конденсат: \textbf{да, $I_B=3$};
\item унаследованный портал $I_B\lambda$: \textbf{нет};
\item лучший следующий кандидат: \textbf{$I_B=\Tr(BB^*)$};
\item следующий гейт: \textbf{parent-origin портала радиуса когерентности}.
\end{enumerate}\end{protocol}

Машинный сертификат сохранён в
\path{s2t_v8_baryon_c0_singlet_triplet_central_gap_existing_scalar_source_carrier_classification_gate_results.json}.

\begin{thebibliography}{9}
\bibitem{v8-existing-scalar-patt-wilczek}
B.~Patt, F.~Wilczek,
\emph{Higgs-field Portal into Hidden Sectors},
arXiv:hep-ph/0605188.
\bibitem{v8-existing-scalar-devastato-lizzi-martinetti}
A.~Devastato, F.~Lizzi, P.~Martinetti,
\emph{Grand Symmetry, Spectral Action, and the Higgs mass},
arXiv:1304.0415.
\bibitem{v8-existing-scalar-stephan}
C.~A.~Stephan,
\emph{New Scalar Fields in Noncommutative Geometry},
arXiv:0901.4676.
\end{thebibliography}